\documentclass[11pt]{article}
\usepackage[a4paper,margin=1in]{geometry}
\usepackage{amsmath,amssymb,amsthm,microtype}
\newtheorem{theorem}{Theorem}[section]
\newtheorem{corollary}[theorem]{Corollary}
\theoremstyle{remark}\newtheorem{remark}[theorem]{Remark}
\title{Finite-Moment Nonidentifiability by Hidden Root-of-Unity Shells}
\author{Bin Wang}\date{July 2026}
\begin{document}\maketitle
\begin{abstract}
Growing coefficient prefixes do not by themselves identify a spectral cloud.
For every prescribed prefix length $N$, we construct a conjugation-symmetric
shell whose power sums vanish through order $N$ but whose next selected
moment is nonzero.  Appending the shell to any spectral multiset preserves
the entire observed prefix while multiplying the regularized determinant by
$1-(\gamma z)^L$ for some $L>N$.  Thus arbitrarily long finite agreement is
compatible with a new divisor immediately beyond the checked window.  The
result is scoped: weighted tail bounds, root transport, or contour control can
exclude the construction.  It shows exactly why the rate-free prefix theorem
of RH-287 cannot replace the weighted gluing leaf of RH-288.
\end{abstract}

\section{Hidden shell}
Fix $L\ge2$ and $\gamma>0$.  Let
\[
 \mathcal S_{L,\gamma}
 =\{\gamma e^{2\pi i j/L}:0\le j<L\}.
\]
This multiset is closed under conjugation.

\begin{theorem}[prefix invisibility]\label{thm:shell}
For every integer $n\ge1$,
\[
 \sum_{\zeta\in\mathcal S_{L,\gamma}}\zeta^n
 =L\gamma^n\mathbf1_{L\mid n}.
\]
Consequently, if $L>N$, every power sum through order $N$ vanishes, while
the order-$L$ sum equals $L\gamma^L$.
\end{theorem}
\begin{proof}
Factor out $\gamma^n$ and sum the finite geometric progression
$\sum_{j=0}^{L-1}e^{2\pi i jn/L}$.
\end{proof}

\begin{corollary}[exact hidden divisor]
The genus-one shell factor is
\[
 \prod_{\zeta\in\mathcal S_{L,\gamma}}
 (1-z\zeta)e^{z\zeta}=1-(\gamma z)^L.
\]
\end{corollary}
\begin{proof}
The linear exponential terms cancel because the first shell moment is zero.
The remaining polynomial identity is the factorization of
$1-(\gamma z)^L$ over the $L$th roots of unity.
\end{proof}

\section{Nonidentifiability consequence}
Let $C$ be any finite spectral multiset and set
$\widetilde C=C\sqcup\mathcal S_{L,\gamma}$.  If $L>N$, then
\[
 \sum_{\lambda\in C}\lambda^n
 =\sum_{\lambda\in\widetilde C}\lambda^n,
 \qquad1\le n\le N,
\]
but the associated canonical products differ by $1-(\gamma z)^L$.
No rule reading only the first $N$ moments can distinguish them.

The same construction can be applied at a moving $N=N_\sigma$.  Without a
weighted estimate that suppresses $\gamma_\sigma^{L_\sigma}R^{L_\sigma}$,
even a diverging prefix does not determine the analytic factor on $|z|\le R$.

\begin{remark}
This counterexample does not show that the physical folded Gaussian spectrum
contains hidden shells.  It only proves insufficiency of prefix data.  A
root-$\ell^1$, weighted-Fourier, or contour theorem would add information that
the construction ignores.  Gates A--E remain false/open.
\end{remark}
\end{document}
